% Plan v5 / item A2-A4: proofs migrated out of the main text, plus the no-free-lunch proof.
% The two extra-pair head-to-head tables were cut on 2026-09-19 for the page budget; both arms are
% unchanged and their CSVs are in results/, and the readings they carried are stated in the prose.
\section{Proofs, and the measurements that check them}\label{app:proofs}

\subsection{Proposition~\ref{prop:threshold}}
Nothing below uses the form of $q$, only the budget it holds to, which is why the statement covers a
metered decoder and a selection rule alike. For $\alpha = \infty$ the budget gives
$q(E) \le e^{K} p_s(E) = e^{K - S(x)}$, which is at
least $1$ exactly when $K \ge S(x)$. For $\alpha \in (1,\infty)$, R\'enyi-to-probability conversion
gives $q(E) \le \left(e^{K} p_s(E)\right)^{(\alpha-1)/\alpha}$; since the exponent is
positive and finite, this is at least $1$ exactly when $e^{K} p_s(E) \ge 1$, i.e. $K \ge S(x)$. For
$\alpha = 1$ the tightest bound is the binary-KL inversion
$\max\{p : d(p \,\|\, e^{-S(x)}) \le K\}$, which attains $1$ exactly when $d(1 \| e^{-S(x)}) = S(x)
\le K$. In all three cases the threshold is $K = S(x)$. \qed

\paragraph{The threshold for a \emph{window} event, which is what extraction actually measures}
\label{app:windowvacuity}
Proposition~\ref{prop:threshold} is stated for an arbitrary event $E$, so it applies verbatim to
$E_w$ = ``reproduce some $w$-token window of the work''. One report asks for that analogue, and it
is sharper than the whole-work statement in the direction that matters. A shorter event is less
surprising under the anchor, so $S_w(x) < S(x)$ and the certificate goes vacuous \emph{sooner};
quoting only the whole-work figure is the lenient reading. It is also the wrong one here, because
\textbf{every extraction number in this paper is a window event} --- near-verbatim recall is scored
on $50$-token windows.

With $S_w = w\,s$ and $s$ the anchor's per-token surprisal on protected text, already measured per
passage by \texttt{analysis/\allowbreak regimes.py} ($758$ CopyBench passages, $16$ novels, median
$3.197$ nats per token, $[2.609, 3.851]$ at the deciles), Table~\ref{tab:windowvacuity} gives the
thresholds and the budgets that reach them at $T_{\max}=200$:

\begin{table}[ht]
\centering\footnotesize
\caption{Vacuity thresholds by event size, from the audited anchor's own per-token surprisal. The
budget column is the $k$ at which $K = k\,T_{\max}$ reaches the threshold at $T_{\max}=200$.}
\label{tab:windowvacuity}
\begin{tabular}{@{}lrrr@{}}
\toprule
event $E$ & $S(x)$, nats & vacuous at $k$ & $\log 64$ sits below by \\
\midrule
$10$-token window   & $32.0$  & $0.160$ & $7.7\times$ \\
$20$-token window   & $63.9$  & $0.320$ & $15.4\times$ \\
$50$-token window   & $159.8$ & $0.799$ & $38.4\times$ \\
whole passage (median, $64$ tokens) & $204.8$ & $1.024$ & $49.2\times$ \\
$100$-token window  & $319.7$ & $1.598$ & $76.9\times$ \\
\bottomrule
\end{tabular}
\end{table}

Two things follow. First, \textbf{at every budget the mechanism's own evaluation uses} --- $k = 1, 3, 5,
10, 20$ --- the certificate is vacuous for a $10$-token window, a $20$-token window, a $50$-token
window and the whole passage alike; the $50$-token event is already vacuous below $k = 0.8$, and
the authors' book setting is $k = 3$. Second, selection's certificate of $\log 64 = 4.159$ nats
--- a bound, against thresholds derived from measured surprisal --- sits below \emph{every}
threshold on this grid, by $7.7\times$ at the smallest window here. The asymmetry
the paper is about does not depend on choosing a generous event, nor on a generous corpus
(Appendix~\ref{app:bookmia} runs the vacuity statement over a hundred books): it survives
shrinking the event
until the budget is the one thing left that grows with the work
(\texttt{results/\allowbreak window\_\allowbreak vacuity.csv}).

% fig:orders (order_invariance.pdf) cut 2026-09-19 for the page budget: it drew the four orders
% against one work, which is exactly what the proof above states, and moved no claim.

\subsection{Theorem~\ref{thm:nfl}}
Write $q$ for the law of the decoder's output sequence under the budget and $p_s$ for the safe
model's. The chain rule for relative entropy gives
$D_{\mathrm{KL}}(q \| p_s) = \mathbb{E}_q\bigl[\sum_t D_{\mathrm{KL}}(q_t \| p_{s,t})\bigr]$, where
$q_t$ is the distribution actually served at step $t$ and $p_{s,t}$ the safe model's at the same
context. For the audited decoder $q_t = p_{\theta_t}$ and the step charge is exactly
$c_1(\theta_t) = D_{\mathrm{KL}}(p_{\theta_t} \| p_{s,t})$, so a per-trajectory budget
$\sum_t c_t \le K$ gives $D_{\mathrm{KL}}(q \| p_s) \le K$. For $\alpha > 1$ the charge dominates
the KL, since $D_\alpha$ is nondecreasing in $\alpha$ \citep{vanerven2014renyi}, so the same
inequality holds for every order in the class.

That chain-rule step is an assumption about the implementation as much as about the mathematics: it
requires the logged sequence spend to be exactly the sum of the per-step KL charges. We checked it
rather than assumed it. Over $1{,}200$ trajectories and $240{,}000$ decode steps at $k=3$, the
largest discrepancy between $\sum_t a_t$ and the recorded total spend is $1.5\times10^{-4}$ nats,
consistent with floating-point accumulation, and the per-step charge agrees exactly with an
independent recomputation from the two logit vectors at every one of those steps.

Donsker--Varadhan \citep{donsker1975asymptotic} states that for every bounded measurable $U$ and
every $\lambda \in \mathbb{R}$,
\[
D_{\mathrm{KL}}(q \| p_s) \;\ge\; \lambda\, \mathbb{E}_q[U] - \log \mathbb{E}_{p_s}\!\left[e^{\lambda U}\right].
\]
Taking the supremum over $\lambda$ gives the convex conjugate of the log-moment generating function
of $U$ under $p_s$, that is the Cram\'er rate function \citep{dembo1998large},
$D_{\mathrm{KL}}(q \| p_s) \ge \Lambda^*_s\bigl(\mathbb{E}_q[U]\bigr)$. Chaining the two displays
gives $K \ge \Lambda^*_s(u)$ with $u = \mathbb{E}_q[U]$, which is the first claim. Since
$\Lambda^*_s$ is convex, vanishes at $u = \mathbb{E}_{p_s}[U]$ and is strictly positive to the right
of it, any $u$ strictly above the safe model's own mean utility forces $K > 0$ and, by
monotonicity, forces $K \ge \Lambda^*_s(u)$ specifically. Proposition~\ref{prop:threshold} then
applies verbatim to any work with $S(x) \le \Lambda^*_s(u) \le K$. \qed

\subsection{Proposition~\ref{prop:outrun}}
\begin{proposition}[Protection outruns certification]\label{prop:outrun}
Under the banking rule, if the decoder emits each token of $x$ with probability at least $1-\eta$
then $k \ge (1-\eta)\,k_{\mathrm{crit}}(x) - \log 2$, and $k_{\mathrm{crit}}(x) \ge s(x) + \delta/N$
for every work, with equality iff the surprisal accumulates at a constant rate.
\end{proposition}

\emph{Proof.} Apply the data-processing inequality to the indicator of the true token: the step
charge is at least the binary divergence between the served and reference probabilities of that
token, and $d(q \,\|\, e^{-s_t}) \ge (1-\eta)s_t - \log 2$ whenever $q \ge 1-\eta$. The banking rule
enforces $\sum_{i \le t} c_i \le (t+1)k - \delta$ at every step, and combining the two for every $t$
gives the first claim. The second, $k_{\mathrm{crit}}(x) \ge s(x) + \delta/N$, is the statement that
a running maximum dominates its own endpoint. \qed

\textbf{The equality condition is not a constant rate, and an earlier version of this appendix said
it was.} Equality holds exactly when the running ratio
$f(t) = (\sum_{i \le t} s_i + \delta)/N_t$ attains its maximum at the last step. At $\delta = 0$ a
constant rate does that, with $f \equiv s(x)$. At $\delta > 0$ it does the opposite: a constant rate
gives $f(t) = s(x) + \delta/N_t$, which is strictly \emph{decreasing}, so the maximum sits at $t=1$
and the inequality is strict. Equality under a positive opening debt needs the surprisal
back-loaded enough to outrun the decaying $\delta/N_t$ --- a work that begins predictably and ends
in its rarest material. Only the data-processing step is specific to this decoder; the charge need
only dominate the KL, which every order $\alpha \ge 1$ does.

$k_{\mathrm{crit}}$ is not our invention: it is the Loynes running maximum representing the workload
of a queue \citep{loynes1962stability}, equivalently the smallest rate of a $(\sigma,\rho)$ arrival
curve admitting the surprisal process \citep{cruz1991calculus,leboudec2001network}. Naming it makes
the diagnosis concrete: a budget publishes a \emph{drift} where safety is a \emph{workload maximum}.
It bounds regime boundaries and is not a per-work screen --- as a predictor of the onset it misses by
$21.9\%$ against a no-change null's $12.3\%$ and ranks the seven pairs it was specified at
$\rho = +0.036$, and all nine at $-0.050$ (Appendix~\ref{app:limitations}).

\subsection{The charge of the bucket-metered class}
Write $\ell = \log p_r - \log p_s$ and $Z(u) = \mathbb{E}_{p_s}[e^{u\ell}]$ along the geometric path
$p_\theta \propto p_s^{1-\theta}p_r^{\theta}$. The charge at order $\alpha$ is
\begin{equation}
c_\alpha(\theta) \;=\; D_\alpha(p_\theta \,\|\, p_s)
 \;=\; \frac{\log Z(\alpha\theta) - \alpha \log Z(\theta)}{\alpha - 1},
\label{eq:renyi-path}
\end{equation}
with $c_1(\theta) = \theta Z'(\theta)/Z(\theta) - \log Z(\theta)$ in the limit, the charge the
deployed decoder uses, and $c_\infty(\theta) = \theta \max_y \ell(y) - \log Z(\theta)$ at the other
end, the pathwise charge Proposition~\ref{prop:selection} gives selection anchoring for nothing.
Theorem~\ref{thm:nfl} needs none of this --- it holds for any law inside the budget --- but
Proposition~\ref{prop:imitation} and every measured arm in Section~\ref{sec:orders} are about this
class specifically.

\subsection{Proposition~\ref{prop:imitation}}
Write $c(\theta) = D_{\mathrm{KL}}(p_\theta \| p_s)$ for the charge at step $t$, dropping the index.
$c$ is nondecreasing on $[0,1]$ with $c(0) = 0$, and the decoder serves
$\theta_t = \max\{\theta \in [0,1] : c(\theta) \le k_t\}$. At a step where the constraint is slack
--- that is, where $c(1) \le k_t$ --- that maximum is attained at $\theta_t = 1$, so the decoder
serves $p_{r,t}$ exactly and is charged $c(1) = D_{\mathrm{KL}}(p_{r,t} \| p_{s,t})$. Writing
$\mathcal{S}$ for the set of slack steps and $\beta = 1 - |\mathcal{S}|/T$,
\[
Z_T \;=\; \sum_{t \le T} c(\theta_t) \;\ge\; \sum_{t \in \mathcal{S}} D_{\mathrm{KL}}(p_{r,t} \| p_{s,t})
\;\ge\; (1-\beta)\, T \cdot \min_{t \in \mathcal{S}} D_{\mathrm{KL}}(p_{r,t} \| p_{s,t}),
\]
the first inequality because every term of the sum is nonnegative. Nothing about the utility enters.

For the second part, $\Lambda^*_s$ is convex, vanishes at $u_0 = \mathbb{E}_{p_s}[U]$ and is
nondecreasing on $[u_0, u_{\max}]$, so $\Lambda^*_s(\mathbb{E}_q[U]) \le \Lambda^*_s(u_{\max})$ for
any $q$; and $\Lambda^*_s(u_{\max}) = -\log \Pr_{p_s}[U = u_{\max}]$. Dividing gives the stated
$\Omega(T)$ provided $\beta$ is bounded away from $1$, the slack steps carry divergence bounded away
from $0$, and that last quantity is $O(1)$ in $T$.

\textbf{The third proviso is not free and the proposition now carries it.} $\Pr_{p_s}[U = u_{\max}]$
is a constant in $T$ only for a utility whose best outcome the safe model reaches at a rate that
does not vanish as the sequence grows. It fails for a utility that only one exact long string
attains, where $\Pr_{p_s}[U = u_{\max}] \approx e^{-cT}$ makes $\Lambda^*_s(u_{\max}) = \Theta(T)$
and the overhead ratio $\Theta(1)$ --- so the separation is a claim about utilities a deployer can
ask for, not about every bounded functional of a sequence. The utility measured in this paper is of
the first kind: winning a judged comparison against the risky model, which the anchor's own
unassisted draw already wins on about a third of prompts, priced at $1.30$ nats below and unchanged
by length.

\subsection{What the propositions are worth in numbers}
Everything from here is measurement rather than proof: the two factors
Proposition~\ref{prop:imitation} multiplies, the causal placement
Proposition~\ref{prop:sparse} permits, and what Theorem~\ref{thm:nfl}'s frontier costs at the
budgets actually served.

\begin{table}[h]
\centering
\small
\caption{\textbf{Both factors of Proposition~\ref{prop:imitation}, measured over $16{,}500$ logged
trajectories at the audited pair.} $\beta$ is the fraction of steps at which the bucket binds, the
imitation rate is the mean $D_{\mathrm{KL}}(p_{r,t}\|p_{s,t})$ over the slack steps the bound is
taken over, and the realised rate is the spend the decoder actually incurs per token. Spend is a
realisation throughout this table; $K = kT_{\max}$ at $T_{\max} = 200$ is the published cap. The
last column regresses cumulative spend on the step index within a trajectory.}
\label{tab:imitation}
\begin{tabular}{rrrrrrrr}
\toprule
$k$ & $K$ & $\beta$ & imitation rate & realised rate & spend & spend\,/\,$K$ & $R^2$ \\
 & (nats) & (binding) & (nats/token) & (nats/token) & (nats) & & (cum.\ vs.\ $t$) \\
\midrule
$0.1$ & $20$ & $0.9540$ & $0.034$ & $0.051$ & $9.3$ & $46.5\%$ & $0.9794$ \\
$0.15$ & $30$ & $0.9249$ & $0.063$ & $0.093$ & $16.9$ & $56.3\%$ & $0.9942$ \\
$0.25$ & $50$ & $0.8507$ & $0.135$ & $0.180$ & $33.2$ & $66.4\%$ & $0.9987$ \\
$0.5$ & $100$ & $0.5915$ & $0.359$ & $0.404$ & $78.5$ & $78.5\%$ & $0.9994$ \\
$1$ & $200$ & $0.2057$ & $0.665$ & $0.676$ & $134.1$ & $67.1\%$ & $0.9940$ \\
$3$ & $600$ & $0.0332$ & $0.828$ & $0.828$ & $165.0$ & $27.5\%$ & $0.9915$ \\
$5$ & $1000$ & $0.0074$ & $0.848$ & $0.850$ & $169.9$ & $17.0\%$ & $0.9915$ \\
$10$ & $2000$ & $0.0001$ & $0.857$ & $0.857$ & $171.3$ & $8.6\%$ & $0.9915$ \\
$20$ & $4000$ & $0.0000$ & $0.857$ & $0.857$ & $171.3$ & $4.3\%$ & $0.9916$ \\
\bottomrule
\end{tabular}
\end{table}

Table~\ref{tab:imitation} is those two factors at the audited pair. Two regimes, in one column. Below the imitation rate the meter binds almost everywhere and the
decoder spends most of its allowance; above it the meter goes slack, the realised rate stops at
$0.857$ nats per token, and the rest is unreachable --- the certificate at $k=20$ is written against
a quantity the decoder uses $4.3\%$ of. The crossover is not a risk level: it is the pair's mean
$D_{\mathrm{KL}}(p_{r,t}\|p_{s,t})$, which is exactly what Eq.~\eqref{eq:req} says a decoder must
\emph{afford}. That half of Eq.~\eqref{eq:req} is confirmed here directly, even though its
prediction of where leakage \emph{begins} was refuted (Appendix~\ref{app:onset}). Here $\beta$
counts every step the bucket touches, including the prefix-debt opening where the anchor writes the
token outright; the strictly interior blends alone are $0.9\%$ of steps at $k=3$. On the $3{,}300$
protected-passage trajectories the same rate is higher, $0.911$ nats per token: the meter charges most
where the risky model is most distinctive from the anchor, which is where memorisation lives. The
last column's variation is across steps within a trajectory, not across trajectories, every one of
which runs to $T_{\max}$. Where the spend goes is the second shape Proposition~\ref{prop:sparse} constrains: the busiest
$1\%$ of steps carry $5.8\%$ of it and covering $90\%$ takes $64\%$ of the sequence, so the deployed
rule sits only modestly off the uniform diagonal and nowhere near the left axis a bounded-budget
policy would need, and once the meter stops binding ($k=3$ and $k=20$ coincide) the shape no longer
depends on the cap at all. At $k=20$ the decoder serves $p_{r,t}$ unchanged at every step it takes
($\beta = 0.0000$), each charged the full $D_{\mathrm{KL}}(p_{r,t}\|p_{s,t})$, so
$N_\varepsilon = \Theta(T)$ for every $\varepsilon$ below the pair's imitation rate --- the vacuous
horn in Proposition~\ref{prop:sparse}'s own currency.

\paragraph{Selection anchoring does not evade the chain rule, and saying so was an error}
Earlier versions of this appendix argued that a selection rule ``is not a causal policy at all'',
that every token it emits is drawn from $p_{s,t}$ so $N_\varepsilon = 0$ identically, and that its
$\log n$ is spent where the chain rule does not reach. The first clause is a statement about the
\emph{sampler} and the rest does not follow from it. The served law $q$ is a distribution over
sequences like any other, so it factorises as $\prod_t q_t(y_t \mid y_{<t})$ and
Proposition~\ref{prop:sparse} applies to it with $K$ its own KL, $3.175$ nats at $n=64$; its
conditionals genuinely differ from $p_{s,t}$, because conditioning on which draw won reweights every
step. Selection \emph{satisfies} the proposition rather than escaping it --- at most $3.175/\varepsilon$
steps above $\varepsilon$, so fewer than four of a $204$-token completion carry a whole nat --- and it
gains $+0.1045$ out of that budget, which is why we now read the second horn as a shape and not a
verdict.

What is true, and is the operative difference, is that those conditionals are not \emph{causally
computable}: $q_t(\cdot \mid y_{<t})$ depends on the scores of \emph{complete} candidates, so no
rule that must commit token $t$ before token $t+1$ exists can serve them. The budget is concentrated
on a decision the decode loop does not contain. That is why the construction below takes the shape
it does: it is the closest a policy that must spend online can come to spending late.

\noindent\emph{Proposition~\ref{prop:sparse}, restated.} Let $q$ be any causal policy,
$q(y) = \prod_{t \le T} q_t(y_t \mid y_{<t})$, with sequence budget
$D_{\mathrm{KL}}(q \,\|\, p_s) \le K$, and let
$N_\varepsilon = \#\{t \le T : D_{\mathrm{KL}}(q_t(\cdot \mid y_{<t}) \,\|\, p_{s,t}) > \varepsilon\}$.
Then $\mathbb{E}_q[N_\varepsilon] \le K/\varepsilon$ for every $\varepsilon > 0$.

\emph{Proof.} The chain rule for an autoregressive factorisation is exact:
$D_{\mathrm{KL}}(q \| p_s) = \sum_{t \le T} \mathbb{E}_{y_{<t} \sim q}\!\left[D_{\mathrm{KL}}(q_t \| p_{s,t})\right]$,
every term nonnegative. Markov's inequality at each step gives
$\Pr_q[D_{\mathrm{KL}}(q_t \| p_{s,t}) > \varepsilon] \le \mathbb{E}_q[D_{\mathrm{KL}}(q_t \| p_{s,t})]/\varepsilon$,
and summing over $t$ bounds $\mathbb{E}_q[N_\varepsilon]$ by $K/\varepsilon$. \qed

\paragraph{Proposition~\ref{prop:sparse} holds surely, not just in expectation, for the deployed class}
The proposition bounds a count in expectation, because $D_{\mathrm{KL}}(q\|p_s)$ is an expectation
over trajectories. The deployed class enforces more: the token bucket caps the \emph{realised}
charges of each trajectory, $\sum_{t \le T} a_t \le K$ pathwise, the accounting under which no
violation occurs in more than $100{,}000$ trajectories. Under it the count is deterministic --- at
most $K/\varepsilon$ steps of any single served trajectory can carry a charge above $\varepsilon$,
since $K \ge \sum_t a_t \ge \varepsilon N_\varepsilon$ --- so the sparsity horn constrains
\emph{every} output the mechanism serves and not merely the average one. It also shows where that horn has no force: at the
authors' own $k=3$ with $T_{\max}=200$ the count is $600/\varepsilon$ against a trajectory of $200$
steps, so it constrains nothing below $\varepsilon = 3$ nats in a single step, which is the first
horn again in another currency.

\paragraph{The dichotomy at finite $T$, and in probability rather than in expectation}
Two reports ask for the dichotomy without an asymptotic in it. Proposition~\ref{prop:sparse} is
already non-asymptotic --- it holds at every $T$ --- but it bounds a mean, and the statement worth
having bounds a tail. Markov applied once more to the count gives, for every $\varepsilon > 0$ and
every $\delta \in (0,1)$,
\[
  \Pr\nolimits_q\!\left[\,N_\varepsilon \ \ge\ \frac{K}{\varepsilon\delta}\,\right] \ \le\ \delta ,
\]
so with probability at least $1-\delta$ a served trajectory carries more than $\varepsilon$ nats at
fewer than $K/(\varepsilon\delta)$ of its $T$ steps. The bound says something only when
$K/(\varepsilon\delta) < T$, that is when
\[
  \varepsilon \ >\ \frac{K}{\delta\,T}.
\]
Both horns are now visible without a limit. If the budget is a \emph{rate}, $K = kT$, the threshold
is $\varepsilon > k/\delta$ --- \textbf{independent of $T$}. Growing the workload never buys
sparsity: at the authors' $k = 3$ and $\delta = 0.05$ the guarantee is silent below $60$ nats in a
single step, at $T = 200$ and at $T = 10^{6}$ alike, and a single step cannot carry $60$ nats
against any anchor that puts non-negligible mass on the served token. If instead the budget is
\emph{bounded}, $K = O(1)$, the same threshold is $\varepsilon > K/(\delta T) \to 0$: for any fixed
$\varepsilon$ and any fixed $\delta$, all but $K/(\varepsilon\delta) = O(1)$ of the $T$ steps are
within $\varepsilon$ of the anchor with probability $1-\delta$, and the decoder \emph{is} the
anchor at all but a vanishing fraction of the work. The choice between a threshold that does not
move with $T$ and one that vanishes with $T$ is the whole of the dichotomy, and neither line of it
needed a limit to state.

\noindent Two further pairs test whether the reversal is the mechanism's or the pair's. Both run
against the same risky model on the same $500$ ordinary prompts, one pass, one shared control, bands
fixed before the selection arm was generated; neither anchor is a \emph{safe} model, so no
certificate, leakage or $s(x)$ number is taken from either.

\emph{Llama-3.2-1B.} The metered decoder is genuinely useful here, resolving at three of four
budgets --- $+0.051$, $+0.056$ and $+0.055$ under judge~B at $k=1,3,20$, against $+0.015$ at
$k=0.5$, for realised spends of $88.1$, $112.7$, $123.2$ and $123.4$ nats --- because a $1$B anchor sits close enough to the risky model for a
per-token meter to buy something. Selection at $n=8$ gains \emph{more}, $+0.076$ $[+0.031, +0.122]$ and $+0.154$
$[+0.104, +0.200]$ on the two judges, for $1.204$ nats: $102\times$ and $94\times$ less divergence,
where the band fixed in advance asked only that it come within $0.03$ at under a tenth of the spend.
The metered spend saturates exactly as Proposition~\ref{prop:imitation} says --- $k=3$ and $k=20$
differ by $0.2$ nats and $0.001$ of gain while the published cap between them differs sevenfold.

The metered decoder fuses two distributions over one vocabulary, so a pair needs a shared vocabulary
between anchor and risky model, and no other openly licensed anchor has one with this risky model.
\emph{Llama-3.2-3B-Instruct}\label{app:frontier3} is the closest anchor to the risky model we can
meter at all, and the protocol named it the hardest case for the reversal, on the theory that a close
anchor is where a per-token meter has most to buy. The prediction was right that the pair is
different and wrong about how. \textbf{The metered decoder resolves at no budget on either judge}:
all eight intervals contain zero and all eight point estimates are negative, the gains running
$-0.038$ to $-0.019$ under judge~B and $-0.030$ to $-0.009$ under judge~C for realised spends of
$54.8$, $59.3$, $60.3$ and $60.6$ nats. Forty times the
published cap moves the realised spend by $10.5\%$ ($54.8$ to $60.6$ nats) and the judged gain by
less than one interval width. A close anchor does not narrow the gap by lifting the meter; it widens
it by flattening it onto the anchor: sparsity that really is triviality, observed here rather
than inferred from Proposition~\ref{prop:sparse}, which does not by itself imply it. Selection over the same anchor gains $+0.150$ $[+0.106, +0.194]$ and $+0.273$
$[+0.227, +0.323]$ for $1.204$ nats, $50\times$ and $49\times$ less than the best metered arm spends
to gain nothing.

\emph{A floor, sighted three times.} Each pass contains two arms that ought to be identical --- the
metered run's own $k=0$ arm and selection's $n=1$ first draw, each one unconstrained completion from
the same anchor on the same prompts from two independent generation runs. Judged, they differ by
$-0.034$ $[-0.071, +0.004]$ and $+0.011$ $[-0.025, +0.048]$ at the first pair and $-0.009$
$[-0.053, +0.037]$ and $+0.010$ $[-0.044, +0.062]$ at the second, no interval excluding zero. The
Comma-7B sweep to $n=64$ gives the same size independently: re-generating the configuration the
breadth arm had run, its $n=8$ gain came back at $+0.072$ against $+0.111$ --- $0.039$, with mean
completion length $80.4$ words against $99.5$ --- and its own nested check declared the cross-arm
reading unreadable. Those are this paper's floor on \emph{generation-run} variation.
\textbf{The cross-pass floor on a judged gain is therefore about $0.04$}, it is
a different quantity from position bias, and every comparison here between gains from
\emph{different} passes carries it. The ones that do not are read inside one pass: the rows above,
and every $n$-sweep.
\paragraph{Proposition~\ref{prop:imitation}'s shape is the mechanism's, not the pair's}
Two further anchors and a $70$B risky model were run on the same $500$ ordinary prompts, budget grid
and $200$-token cap, with both mandatory baselines and bands fixed in advance. They vary the distance
between the two distributions while the mechanism, the workload and the judges stay fixed.

\begin{table}[h]
\centering\footnotesize
\caption{\textbf{Proposition~\ref{prop:imitation}'s shape at four pairs.} $r_{\rm imit}(k)$ is the
imitation rate measured at budget $k$; the ratio column tests saturation against a band of $1.10$
fixed in advance, and $|r_{\rm real}/r_{\rm imit}-1|$ how closely realised spend converges on it.
The level moves with how far the risky model has been tuned from the anchor, not with its parameter
count.}
\label{tab:imitation-pairs}
\setlength{\tabcolsep}{4pt}
\begin{tabular}{@{}lccccc@{}}
\toprule
anchor & $r_{\rm imit}(3)$ & $r_{\rm imit}(20)$ & ratio & $|r_{\rm real}/r_{\rm imit}-1|$
& spend, $k{=}20$ \\
\midrule
TinyComma-1.8B, 8B-Inst.  & $0.8278$ & $0.8570$ & $1.035$ & $0.0005$ & $171.3$ \\
Llama-3.2-1B, 8B-Inst.    & $0.6157$ & $0.6174$ & $1.003$ & $0.0005$ & $123.4$ \\
Llama-3.2-3B-Inst., 8B-Inst. & $0.3010$ & $0.3034$ & $1.008$ & $0.0013$ & $\ \ 60.6$ \\
TinyComma-1.8B, 70B base  & $0.8330$ & $0.9085$ & $1.091$ & $0.0007$ & $148.9$ \\
\bottomrule
\end{tabular}
\end{table}

All four readings hold at all four pairs: the rate saturates (the ratio column, against a specified
$1.10$), the realised spend converges on it to within $0.14\%$, the cumulative spend is linear in the
step index at every budget with a minimum median $R^2$ of $0.974$, and at $k=20$ the decoder serves
$p_r$ unchanged at $\ge 99.95\%$ of steps. What moves is the \emph{level}, and the fourth pair says
which distance it tracks. The specified band read a $70$B base risky model as \emph{further} from a
$1.8$B anchor than an $8$B instruction-tuned one and predicted a higher rate; it came back flat,
$0.8330$ against $0.8278$, and an addendum specified before the arm finished had already named the
confound --- the $8$B is instruction-tuned and the anchor is not, so the $0.8278$ already contains a
base-to-instruct shift the $70$B does not pay. Read as \emph{tuning} rather than parameter count the
levels order cleanly, $0.303$, $0.617$, $0.857$, $0.908$, with realised spend following exactly. The
prediction that the same-family instruct rate would fall below half the same-family base rate passed narrowly, by $2.2\%$ of its threshold on $r_{\rm imit}(3)$, the column it was specified on, and is
reported as narrow.
\paragraph{The causal placement Proposition~\ref{prop:sparse} permits}
The front-loaded decoder is granted $\log 8 = 2.0794$ nats up front rather than metered per token,
and spends them, realised median \emph{and} maximum $2.0794$. On the same $500$ prompts and judge it
gains $-0.008$ $[-0.056, +0.041]$; the uniform rate at the same total gains $-0.032$
$[-0.083, +0.016]$ and ten times the budget gains $-0.019$. The two placements differ by $0.024$
with overlapping intervals, so the paper claims no ordering between them --- only that neither
resolves, where selection at the same $2.08$ nats gains $+0.054$. All three bands were fixed before
the runs.

A later arm extends this in the two directions that matter, judged through the order-averaged
instrument rather than the placement script, so its numbers are not quoted across the ones above. In
\emph{budget}, the policy gains $-0.0015$ $[-0.0200,+0.0170]$ at $\log 8$ nats, $+0.0215$
$[+0.0015,+0.0415]$ at $\log 64$ and $+0.0545$ $[+0.0315,+0.0775]$ at $64$: the first agrees with the
$-0.008$ above and the other two exclude zero, so \textbf{the causal horn is not empty} --- it is
simply far from the frontier, since at matched budget the paired difference is $+0.083$
$[+0.052,+0.1135]$, no budget on the grid reaches selection's $+0.1045$ at $\log 64$, and at $64$
nats, $15.4\times$ the budget, the interval tops out at $+0.0775$. Nor do the orders favour the causal
arm, since $D_\infty \le \log 64$ implies $D_{\mathrm{KL}} \le \log 64$ while the $64$-nat arm bounds
only $D_{\mathrm{KL}}$. In \emph{placement}, granting a budget up front is not choosing where to spend
it, so the decoder was given a reserving rule: spend at step $t$ only if the full-tilt demand
$D_{\mathrm{KL}}(p_{r,t}\Vert p_{s,t})$ reaches $\tau$ nats, otherwise serve the anchor and keep the
nats. It reserves as designed, the spend moving from median step $3$ to median $17$, and every
threshold loses to it --- $+0.0215$, $+0.0150$, $+0.0110$, $+0.0045$, $+0.0070$ at
$\tau = 0,1,2,4,8$, the last a tick above the fourth rather than a continued fall --- because a higher
bar releases the budget less often ($0\%$ to $86.7\%$ of trajectories never spend) \emph{and} buys
less when it does ($+0.0215$ to $+0.0081$ conditional on spending). The one reading that favours
concentration is reported and not leaned on: at $\tau=8$ the conditional gain is the grid's largest,
$+0.0342$, but on $n=73$ with interval $[-0.0068,+0.0753]$, which includes zero and overlaps
$\tau=0$'s almost entirely.

\paragraph{Proposition~\ref{prop:selection} is tight}
If $f$ is maximised at a single $y^\star$ with $p_s(y^\star) = p$, then
$q(y^\star) = 1-(1-p)^n \to np$ as $p \to 0$, so no constant below $n$ works in
$q(y) \le n\,p_s(y)$ and the $\log n$ charge cannot be improved without assuming something about
$f$ or $p_s$ --- which is the point of assuming neither.

\paragraph{What Theorem~\ref{thm:nfl} is worth in numbers}
The utility is the judge's verdict scored $1/\tfrac12/0$ and its law under $p_s$ is read off the
anchor-only arm; realised spend is the decoder's own sequence divergence, which never reaches its cap.
At $k=0.5,1,3,20$ the decoder spends $79.0$, $134.6$, $165.0$ and $171.3$ nats to reach
$\mathbb{E}_q[U] = 0.334$, $0.420$, $0.469$ and $0.483$, where the rate function $\Lambda^*_s$ of that
utility under the safe model asks for $0.0003$, $0.023$, $0.052$ and $0.061$: every row satisfies
Eq.~\eqref{eq:nfl} by three to five orders of magnitude ($2.8\times10^{5}$, $5900$, $3200$ and $2800$
times over), and bootstrapping over judged pairs at the \emph{weakest} end of each interval still
leaves the ratio between $1.3\times10^{3}$ and $8.6\times10^{3}$. The two mechanisms sit at very
different distances from the frontier under judge~B's own law of $U$: selection at $n=8$ is $58\times$
it, and the metered decoder's best arm $7{,}995\times$.

\paragraph{Why better scheduling cannot close the gap}
The decoder trades two quantities and on the geodesic both are closed forms in $\psi(u) = \log Z(u)$:
\begin{equation}
C(\theta) = \theta\psi'(\theta) - \psi(\theta), \quad C'(\theta) = \theta\psi''(\theta), \qquad
G(\theta) = \theta m - \psi(\theta), \quad G'(\theta) = m - \psi'(\theta),
\label{eq:charge-gain}
\end{equation}
with $m = \psi'(1)$. $C$ is the charge and $G(\theta) = -D_{\mathrm{KL}}(p_r \| p_\theta)$ up to a
per-step constant, so $G$ is fidelity to the model the decoder is imitating, $G(0) = 0$ and
$G(1) = D_{\mathrm{KL}}(p_r\|p_s)$. Note $C'(0) = 0$ while $G'(0) = m - \psi'(0)$ is the Jeffreys
divergence between the two models --- the quantity governing the exact reward--KL
frontier~\citep{monteiropaes2026limits} --- so at every step the first nat bought is worth a Jeffreys
divergence and the last is worth nothing.

The token bucket is myopic: it takes whatever marginal price $G'/C'$ its current allowance lands on.
The best allocation of a \emph{fixed} total instead equalises that price, serving the $\theta_t$ that
solve $G'(\theta) = \lambda C'(\theta)$ at one shared $\lambda$. Both are computable from two
teacher-forced forward passes, with no decoding run and no change to the decoder, so what better
scheduling is worth can be measured rather than argued. We walk both models along $30$ targets ($12$
for KL3M-520M, the one arm run at a smaller limit) and give the equal-price allocation exactly the
total the bucket itself spent, over three kinds of target, because the geometry could in principle
depend on what is being written: a passage the risky model was fine-tuned to reproduce, a novel from
a split it has never seen, and a sample it drew itself on an ordinary prompt.

Optimal reallocation of the same budget buys $3$ to $13\%$ more fidelity at the budgets where the
certificate says anything ($1.123$ and $1.133$ at $k=0.25$ falling to $1.006$ and $1.007$ at $k=3$,
for KL3M-520M and Pleias-1.2B), so a primal-dual scheduler is worth percent, not orders of
magnitude. The three target columns agree to within two points of gain everywhere, so this is a
property of the geometry and not of what is being written. A second reading does not need the first:
the bucket already captures $46.7$--$54.5\%$ of the $\theta = 1$ fidelity at $k=0.25$ and
$92.7$--$98.9\%$ at $k=3$, and $\theta = 1$ everywhere is the risky model served outright, which is
what an unlimited budget buys --- so at $k=1$ the ceiling on \emph{every} allocation at once is only
$1.2$ to $1.3\times$ away, and on the memorised trajectory, the one an extraction adversary walks,
about $1.3\times$. No decoder that still meters divergence from the anchor can recover three orders
of magnitude, because three orders of magnitude of fidelity are not there to recover.

Two limitations. The allocation is computed along a fixed teacher-forced trajectory, and a decoder
that spent differently would walk a different one; and fidelity to $p_r$ is not judged utility. The
second is the point rather than a caveat: the mechanism has no utility signal, and imitating $p_r$ is
the only thing its budget can buy. What this adds to Theorem~\ref{thm:nfl} is that the overhead is
the price of that target, not of the schedule.
% Moved out of Section 4: the enumeration of failed repairs supports the dichotomy,
% it is not the headline, and the main text names all four in prose.
% tab:repairs cut 2026-09-19 for the page budget. Section~\ref{sec:orders} names all four escapes
% and where each lands in prose, and each one's numbers are in the paragraphs above; the table was
% a second telling of both. The row that was only here -- the vacuity count -- is kept below.
\paragraph{Where the four escapes land}
Every escape the dichotomy leaves a per-token meter is measured above: sharpening the charge
($\alpha: 1 \to 8$ at $k=3$) cuts oracle recall $0.0965 \to 0.0012$, an $80\times$ cut at no
measurable utility cost, and lands on the bounded horn because it stops serving the risky model
rather than metering it better; lowering the cap to $k=0.5$ leaves neither outside judge separating it from
its own safe model over $600$ comparisons each ($-1.03\sigma$, $+1.52\sigma$), both separating only at
$k \ge 10$ ($-6.97\sigma$, $-3.29\sigma$); scheduling it better buys $3$ to $13\%$ at $k \le 1$ and
under $1.4\%$ at $k=3$ against an \emph{offline} optimum no causal scheduler can beat
\citep{tomasi2026primaldual}; and granting it up front to spend sparsely reaches $+0.0545$
$[+0.0315, +0.0775]$ at $64$ nats, so the horn is not empty but never reaches selection's $+0.1045$ at
$\log 64 = 4.16$. None of it is visible in the published number: at the authors' own $k=3$ every order
publishes $K = 3T$ and the certificate is vacuous for $100\%$ of the $758$ protected passages and
$98.7\%$ of $9{,}870$ across a hundred books.

